\addtocontents{toc}{\protect\Needspace{3\baselineskip}}
\originalchapter{2007 年第二次正式测验}
原卷为 B. Zwiebach 于 2007 年 5 月 2 日的 Test 2, 限时两小时, 允许个人两页笔记, 末页有公式表. 四题保留所有子问与必要公式, 解答为本稿独立补写. 原卷粒子参数为 $x^+=p^+\tau/m^2$; 若常数重缩放 $\tau$, 补偿参数也须相应转换.

\src{8.251 Test 2 Q1, (a)--(c), 15 分.}
\begin{exercise}
复标量场作用量为 $S=\int\dd^Dx\,\mathcal L$,
$\mathcal L=-\frac12\eta^{\alpha\beta}(\partial_\alpha\phi^*)(\partial_\beta\phi)$.
(a) 若 $\phi\mapsto\ee^{\ii\theta}\phi$, $\theta$ 为实常数, 给 $\phi^*$ 的变换并验证对称性.
(b) 写无穷小变换并将 $\phi,\phi^*$ 视为独立场, 用 $j^\alpha=\sum_A[\partial\mathcal L/\partial(\partial_\alpha\phi^A)](\delta\phi^A/\epsilon)$ 求 Noether 流.
(c) 用 $\phi=A+\ii B$ 两个实场表示所得流.
\end{exercise}
\begin{solution}
$\phi^*\mapsto\ee^{-\ii\theta}\phi^*$, 常相位不被导数作用, 两相位抵消, 故 $\mathcal L$ 不变. 小 $\theta=\epsilon$ 时 $\delta\phi=\ii\epsilon\phi$, $\delta\phi^*=-\ii\epsilon\phi^*$, 得
\[
 j^\alpha=-\tfrac12\partial^\alpha\phi^*(\ii\phi)
 -\tfrac12\partial^\alpha\phi(-\ii\phi^*)
 =\frac\ii2(\phi^*\partial^\alpha\phi-\phi\partial^\alpha\phi^*).
\]
方程为 $\partial^2\phi=\partial^2\phi^*=0$, 求散度得零. 代入 $A+\ii B$, 有
$j^\alpha=B\partial^\alpha A-A\partial^\alpha B$.
也可直接用 $\delta A=-\epsilon B,\delta B=\epsilon A$ 与两个实标量的动量求同一结果. 流的整体正负取决于相位变换约定, 本题已固定为正 $\theta$.
\end{solution}

\src{8.251 Test 2 Q2, (a)--(c), 25 分.}
\begin{exercise}
闭弦光锥规范 $X^+=\alpha'p^+\tau$, $\sigma\in[0,2\pi]$, 取
\[
 X^1=0,\qquad X^2=a\cos\sigma\cos\tau,\qquad
 X^3=a\sin\sigma\cos\tau,\quad a>0,
\]
其余横向方向为零. (a) 给 $x_0^-=0$, 用
$\dot X^-\pm X'^-=(\dot X^I\pm X'^I)^2/(2\alpha'p^+)$ 求 $X^-$.
(b) 由 $X^1=0$ 求 $p^+$.
(c) 求 $M^2$.
\end{exercise}
\begin{solution}
横向速度平方为 $a^2\sin^2\tau$, 切线平方为 $a^2\cos^2\tau$, 内积为零, 所以两支平方均等于 $a^2$,
\[
 X'^-=0,\qquad \dot X^-=\frac{a^2}{2\alpha'p^+},
 \qquad X^-=\frac{a^2}{2\alpha'p^+}\tau.
\]
$X^1$ 是目标方向, 不在横向集合 $I=2,\ldots,D-1$ 内.
由 $X^1=(X^+-X^-)/\sqrt2=0$, $p^+=a/(\sqrt2\alpha')$, 取正能量支. 零模给 $p^-=p^+$, 横向总动量为零, 故
$M^2=2p^+p^-=a^2/\alpha'^2$.
在 $\tau=0$ 的静止圆上, 长度 $2\pi a$, 能量 $T_s2\pi a=a/\alpha'$ 与静止质量相同. 圆会收缩到退化构型, 不能说整个历史每一时刻均为非退化光滑世界面.
\end{solution}

\src{8.251 Test 2 Q3, 三小问合计 20 分.}
\begin{exercise}
(a) 比较 $D$ 维引力子的极化数与 $D-1$ 维一个引力子、Maxwell 场和标量的总极化数.
(b) 在自由粒子光锥量子化中求
$[x^+,p^-]$, $[p^-,x^I]$, $[x_0^-,p^+]$.
(c) 将 $L_{-1}^\perp L_{-1}^\perp|p^+,\boldsymbol p_T\rangle$ 化为正规序振子作用于真空的形式.
\end{exercise}
\begin{solution}
(a) 无质量极化数为 $D(D-3)/2$. 在 $D-1$ 维相应总数为
\[
 \frac{(D-1)(D-4)}2+(D-3)+1=\frac{D(D-3)}2.
\]
按通常传播场解释可取 $D\geq4$. 这是降维后张量分解的自由度检查, 尚不证明全部相互作用匹配.

(b) $x^+=\kappa p^+\tau$（原公式表 $\kappa=1/m^2$）不是独立正则坐标, 而 $p^-$ 是动量函数, 第一式为零. $x^I=x_0^I+\kappa p^I\tau$ 给第二式 $-\ii p^I/p^+$, 第三式为 $-\ii$.

(c) 令 $A=\alpha_{-1}^I\alpha_0^I$. 对任意横向动量真空, $L_{-1}|p\rangle=A|p\rangle$, 且
$[L_{-1},A]=\alpha_{-2}\cdot\alpha_0$, 因 $[L_{-1},\alpha_0]=0$. 因而
\[
 L_{-1}^2|p\rangle=
 [(\alpha_{-1}\cdot\alpha_0)^2+\alpha_{-2}\cdot\alpha_0]|p\rangle,\qquad
 \alpha_0=\sqrt{2\alpha'}p_T.
\]
若额外取零横向动量, 两项都消失.
\end{solution}

\src{8.251 Test 2 Q4, (a)--(e), 40 分.}
\begin{exercise}
临界玻色开弦位于 $x^a=0$ 的 D$p$ 膜. $X^+,X^-$ 沿膜, 横向分为 $i=2,\ldots,p$ 的 NN 方向与 $a=p+1,\ldots,D-1$ 的 DD 方向. 光锥 $X^+=2\alpha'p^+\tau$, $0\leq\sigma\leq\pi$. 已知
\[
 \dot X^-\pm X'^-=\frac{(\dot X^i\pm X'^i)^2+
 (\dot X^a\pm X'^a)^2}{4\alpha'p^+},\quad
 P^{\tau\mu}=\frac{\dot X^\mu}{2\pi\alpha'}.
\]
NN 满足 $\dot X^i\pm X'^i=\sqrt{2\alpha'}\sum_n\alpha_n^i\ee^{-\ii n(\tau\pm\sigma)}$;
延拓流 $A^i(\tau,\sigma)$ 有
$[A^i(\sigma),A^j(\sigma')]=4\pi\ii\alpha'\delta^{ij}\partial_\sigma\delta_{2\pi}(\sigma-\sigma')$.
DD 展开为
\[
 X^a=\sqrt{2\alpha'}\sum_{n\ne0}\frac{\alpha_n^a}{n}
 \ee^{-\ii n\tau}\sin n\sigma.
\]
(a) 求 $X'^a\pm\dot X^a$, 与 NN 比较, 写 $X^a,P^{\tau b}$ 等时对易关系. 定义 $A^a(\tau,\pm\sigma)=X'^a\pm\dot X^a$, 给延拓对易式.
(b) 求新的正规序 $L_0^\perp$ 和质量式, 用 $2\alpha'p^+p^-=L_0^\perp-1$.
(c) 给真空标签, 列 $N=0,1,2$ 全部类型, 解释膜上无质量场.
(d) $\Omega_p$ 同时反转弦定向 $\sigma\mapsto\pi-\sigma$ 与反射法向坐标 $X^a\mapsto-X^a$. 写它对坐标和振子的作用.
(e) 取无额外标签且真空为正, 给投影后 $N=0,1,2$ 的态和 $\Omega_p$ 的层数表达式.
\end{exercise}
\begin{solution}
(a) 求导得
\[
 X'^a\pm\dot X^a=\sqrt{2\alpha'}\sum_{n\ne0}
 \alpha_n^a\ee^{-\ii n(\tau\pm\sigma)}.
\]
与 NN 形式相同, 但没有 DD 零模, 原始流中时间、空间导数的顺序交换. 正则核为
\[
 [X^a(\sigma),P^{\tau b}(\sigma')]=\ii\delta^{ab}\delta_D(\sigma,\sigma'),
 \quad\delta_D=\frac2\pi\sum_{n>0}\sin n\sigma\sin n\sigma',
\]
同时 $[X^a,X^b]=[P^{\tau a},P^{\tau b}]=0$. 在 $[-\pi,\pi]$ 按题意拼成 $A^a$, 得
$[A^a(\sigma),A^b(\sigma')]=4\pi\ii\alpha'\delta^{ab}\partial_\sigma\delta_{2\pi}(\sigma-\sigma')$.
Fourier 投影给 $[\alpha_m^a,\alpha_n^b]=m\delta^{ab}\delta_{m+n,0}$. 端点须保留对应奇延拓, 不把整条直线的核无条件照搬.

(b) 令 $N=\sum_{n>0}(\alpha_{-n}^i\alpha_n^i+\alpha_{-n}^a\alpha_n^a)$, 则
$L_0^\perp=\alpha'p^ip^i+N$, $M^2=(N-1)/\alpha'$.
DD 振子增加层数, 法向连续动量却不存在.

(c) 真空为 $|p^+,p^i\rangle$, 层 0 只此态. 层 1 是
$\alpha_{-1}^i|p\rangle,\alpha_{-1}^a|p\rangle$.
层 2 是 $\alpha_{-2}^I|p\rangle$ 与全部对称
$\alpha_{-1}^{(I}\alpha_{-1}^{J)}|p\rangle$, $I$ 遍历 $i,a$, 再分为 $ii,ia,aa$ 三类. 合计 $24+300=324$. 无质量层的 $p-1$ 个 NN 态为膜上 Maxwell 极化, $D-p-1$ 个 DD 态为法向位置标量.

(d) 有
$\Omega_pX^i(\tau,\sigma)\Omega_p^{-1}=X^i(\tau,\pi-\sigma)$,
$\Omega_pX^a(\tau,\sigma)\Omega_p^{-1}=-X^a(\tau,\pi-\sigma)$.
余弦反射给 $(-1)^n$, 正弦给 $(-1)^{n+1}$ 再乘反射负号, 两类均有
$\Omega_p\alpha_n^I\Omega_p^{-1}=(-1)^n\alpha_n^I$.
零模标签不变.

(e) 正真空给 $\Omega_p=(-1)^N$, 因而层 0 保留, 层 1 全删, 层 2 的全部 324 态保留. 若另有 Chan--Paton 标签, 还须转置并选择相容矩阵表示, 此结论不代表含标签实际 orientifold 模型的全部无质量谱.
\end{solution}

